% ------------------------------------------------------------------------
% file `revisions-exercise-10-exercise-body.tex'
%   in folder `xsim/'
%
%     exercise of type `exercise' with id `10'
%
% generated by the `exercise' environment of the
%   `xsim' package v0.20c (2021/02/03)
% from source `revisions' on 2021/12/14 on line 231
% ------------------------------------------------------------------------
  Résous les problèmes suivants. Fais en sorte que ta réponse soit correctement présentée.
  \begin{enumerate}
    \item Valentine lance une rumeur à 9 h 00. Elle le dit à 3 personnes. Chaque
          heure, toutes les nouvelles personnes au courant le disent à 3
          autres personnes. Combien de personnes seront mises au courant à
          14 h 00~? Et à 17h00~?
    \item L'être humain cligne plus de 10 000 fois par jour des yeux. Combien
          de fois une jeune fille de 20 ans a-t-elle cligné des yeux depuis sa
          naissance~?\par
          Donne le résultat sous la forme scientifique en estimant qu'une année
          compte 365 jours.
    \item Les grains de sable de la plage de Syracuse sont fins puisqu'il en
          faut 10 pour faire un volume de \SI{1}{mm^3}. Il y a du sable sur une
          épaisseur de 1 m, la plage fait 50 m de large sur 2 km de long. Exprime en notation scientifique la quantité de grains de sable présents sur cette plage.
    \item À chaque bond supplémentaire, une grenouille double la longueur de son
          bond. Elle commence par un bond de 2 cm. Quelle distance a-t-elle
          parcourue au total si elle bondit cinq fois~?
    \item Le corps humain renferme environ 5 litres de sang. Sachant qu'il y a 5
          millions de globules rouges dans \SI{1}{mm^3} de sang, écris en notation
          scientifique la quantité de globule rouges qu'un corps humain
          renferme.\\\emph{Note~: \SI{1}{mm^3} = \SI{0,000001}{L}}
    \item Range les planètes suivantes dans l'ordre croissant de masse.
          \begin{center}
            \begin{tblr}{row{1}={font=\bfseries}, hlines, vlines, rowsep=1ex}
              Planète & Masse (kg)\\
              Mercure & \num{33.02e22}\\
              Vénus & \num{4.8685e24}\\
              Terre & \num{597.3e22}\\
              Jupiter & \num{1.8986e27}\\
              Uranus & \num{8.6832e25}\\
              Mars & \num{0.64185e24}
            \end{tblr}
          \end{center}
  \end{enumerate}
